\chapter*{How to Read This Book}
\addcontentsline{toc}{chapter}{How to Read This Book}
\markboth{HOW TO READ THIS BOOK}{HOW TO READ THIS BOOK}

\begin{flushright}
\itshape Lasciate ogne fretta, voi ch'intrate.\\[2pt]
\upshape\small not quite Dante, \emph{Inferno}, canto III
\end{flushright}

Over the gate of Hell, Dante wrote: abandon all hope, ye who enter here. This
book asks for less: abandon all haste.

This is a dense book. It is packed with names, dates, versions and code, and
almost every page leans on something said before it. It assumes you will stop,
go back a few pages, check a date, and come back. It was written for curious
people: the ones who want to know why a thing is the way it is, who would rather
spend an evening on one page than skim twenty, and who are not satisfied until
the idea and the history behind it both make sense.

One project runs through the whole book: the import process of a foreign trade
company. An invoice arrives, someone reads it, someone classifies every item
under a customs code, and the result goes into the company's ERP. When the
company won a large new client, doing that by hand meant doubling the customs
team. The project replaces the inbox with a system: a service that reads
documents, a service that suggests the customs classification, two event buses
and a single master that drives each process from start to end, all on AWS.

Each chapter starts from a decision in that project and asks why the
technologies we reach for are shaped the way they are. The chapters come in
the order the project reaches each subject. None of them replaces the
documentation of the tools it discusses. Each one goes as far as the project
needs and stops there.

\section*{Every idea has a parent}

No technology in this book showed up because it was elegant. The German
industry wrote Profibus in the late 1980s because every machine on a factory
floor was wired to the controller with its own cable, and one shared bus with
one master was cheaper and easier to reason about. Hector Garcia-Molina and
Kenneth Salem described sagas in 1987 because a long transaction that holds its
locks for hours blocks everyone else. Gregor Kiczales and his group at Xerox
PARC named aspect-oriented programming in 1997 because logging, transactions
and security ended up copied into every method of a program. Each idea fixed
something its predecessor got wrong and brought costs of its own. If you know
the lineage, you know when the idea applies. If you only know the rule, you
will apply it where it does not belong.

\section*{Five kinds of boxes}

Stories, concepts, questions, reflections and words are kept apart, so you
always know which one you are reading.

\begin{history}{year \textperiodcentered{} place}
An amber box tells a story: who did it, when, and what forced the idea into
existence. The people, the dates and the problems in it are real. No invented
incidents.
\end{history}

\begin{concept}{name}
A blue box defines an idea you will need again. It always ends the same way:

\textbf{What it fixes:} the problem that made the idea worth having.

\textbf{What it costs:} what you give up by using it.
\end{concept}

\begin{question}{the question you were embarrassed to ask}
A teal box asks the question that occurs to anyone reading the paragraph above
it and that nobody asks out loud, and then answers it.
\end{question}

\begin{reflection}{the thing that should bother you}
A purple box stops on a statement that sounds absurd, or that the text just said
in passing as if it were obvious, and works out what it means and why it is
true.
\end{reflection}

\begin{etymology}
A gray box says where a word comes from. The origin of a word is not its
definition, but it often shows what the people who coined it were looking at.
\end{etymology}

The text between the boxes is the argument that connects them.

\section*{Code}

The examples are in Java and Python, the languages of the services and of the
functions in this project. Source code sits on a light gray background and what
a command printed sits on a dark one. Every example was run before it went into
the book, and the output shown is the output it gave. There is as little code
as possible: a block appears only when it shows something a sentence can't.

\section*{Diagrams}

Colors in diagrams are roles, not decoration. A color means the same thing in
every figure of the book.

\begin{center}
\begin{tikzpicture}
  \matrix[column sep=6mm, row sep=4mm, nodes={anchor=center}] {
    \node[hf/source, minimum width=26mm] {source}; &
    \node[hf/note, anchor=west, text width=95mm, align=left] {where the data comes from: a document, a sender, an outside system}; \\
    \node[hf/reader, minimum width=26mm] {service}; &
    \node[hf/note, anchor=west, text width=95mm, align=left] {a program that does the work: a service, a function, a screen}; \\
    \node[hf/store, minimum width=26mm] {store}; &
    \node[hf/note, anchor=west, text width=95mm, align=left] {where data is kept: a database, a bucket, the ERP}; \\
    \node[hf/bus, minimum width=26mm] {middle}; &
    \node[hf/note, anchor=west, text width=95mm, align=left] {something that sits between: a queue, a bus, a protocol}; \\
    \draw[hf/ask] (-13mm,0) -- (13mm,0); &
    \node[hf/note, anchor=west, text width=95mm, align=left] {asks: a request, a poll}; \\
    \draw[hf/answer] (-13mm,0) -- (13mm,0); &
    \node[hf/note, anchor=west, text width=95mm, align=left] {answers a request}; \\
    \draw[hf/push] (-13mm,0) -- (13mm,0); &
    \node[hf/note, anchor=west, text width=95mm, align=left] {pushes: an event, sent without being asked}; \\
    \draw[hf/save] (-13mm,0) -- (13mm,0); &
    \node[hf/note, anchor=west, text width=95mm, align=left] {keeps: written to storage}; \\
    \draw[hf/pain, dashed] (-13mm,0) -- (13mm,0); &
    \node[hf/note, anchor=west, text width=95mm, align=left, text=hfAlert] {a cost, a pain, or something forbidden}; \\
  };
\end{tikzpicture}
\end{center}

\section*{Exact facts, firm opinions}

This is not a scientific book, and it does not pretend to be one.

The history is checked. When the book says who did something and when, that is
what happened.

The opinions are another matter. When a framework survives only because the
tutorials keep teaching it, the book says so. When a message broker adds nothing
to a system with one reader, the book says it adds nothing, not that it is ``of
limited value in some contexts''. Sometimes that will be unfair or imprecise,
and we accept that.

A reader who leaves with a position can argue with it, refine it, or throw it
away. A reader who leaves with ``it depends'' has nothing to hold. Between being
precise and leaving you with a position, this book prefers to give you a
position.
